\chapter{月相}\label{ch:g7:moon-phases}

几年前你记过一本月亮日记，看着银色的月面形状
从\omterm{ex:g2:moon-in-the-sky:shapes}{蛾眉月}变成\omterm{ex:g2:moon-in-the-sky:shapes}{满月}，再变回去，当时许诺以后再讲“为什么”。
“以后”就是现在。你已经具备解释所需的全部条件：一个
每月绕行一周的\omterm{def:g2:moon-in-the-sky:moon}{月球}、沿直线传播的阳光，以及一个
始终有一半被照亮的球。

\section{背后的一个事实}

\begin{proposition}[始终有一半被照亮]\label{prop:g7:moon-phases:halflit}
太阳在任一时刻都只照亮\omterm{def:g2:moon-in-the-sky:moon}{月球}的一半——面向太阳的那一半，
这与它照亮地球一半的情形完全相同。\omterm{def:g2:moon-in-the-sky:moon}{月球}没有
变化的形状，没有增长或缩小：它是一个石球，
永远有一面向阳、一面背阴，两者的
分界线随它绕行而移动。我们所说的“\omterm{def:g7:moon-phases:phases}{月相}”
不过是\emph{视角}：向阳面中恰有多少
朝向地球。
\end{proposition}

\begin{definition}[月相]\label{def:g7:moon-phases:phases}
\omterm{def:g2:moon-in-the-sky:moon}{月球}的\emph{月相}\index{月相}，是指它在一个月中呈现出的
视觉形状：\emph{新月}（向阳面完全背离我们，什么也看不见）、
\emph{\omterm{ex:g2:moon-in-the-sky:shapes}{蛾眉月}}、\emph{上弦月}
（圆面一半被照亮）、\emph{凸月}（超过一半）、
\emph{\omterm{ex:g2:moon-in-the-sky:shapes}{满月}}（向阳面完全正对我们），再到凸月，
\emph{下弦月}，再回到\omterm{ex:g2:moon-in-the-sky:shapes}{蛾眉月}，完成一个循环所需
约 $29.5$ 天。从新月到\omterm{ex:g2:moon-in-the-sky:shapes}{满月}，被照亮的部分逐夜
增大，此时\omterm{def:g2:moon-in-the-sky:moon}{月球}处于\emph{渐盈}\index{渐盈}状态；
由\omterm{ex:g2:moon-in-the-sky:shapes}{满月}回到新月则逐渐缩小，此时为\emph{渐亏}\index{渐亏}。
\end{definition}

\begin{omfigure}
\begin{tikzpicture}[scale=0.8]
  % dashed orbit
  \draw[dashed, omExo!70] (0,0) ellipse (4.2 and 3.1);
  % eight stations: dark half toward the left, sunlit half (white)
  % toward the sunlight arriving from the right
  \foreach \a in {0,45,...,315} {
    \begin{scope}[shift={({4.2*cos(\a)},{3.1*sin(\a)})}]
      \fill[omExo] (0,0) circle (0.46);
      \begin{scope}
        \clip (0,0) circle (0.46);
        \fill[white] (0,-0.5) rectangle (0.5,0.5);
      \end{scope}
      \draw[thick] (0,0) circle (0.46);
      \draw[thick] (0,0.46) -- (0,-0.46);
    \end{scope}
  }
  % the Earth, its own day side facing the sunlight
  \fill[omDef!35] (0,0) circle (0.55);
  \begin{scope}
    \clip (0,0) circle (0.55);
    \fill[omDef!12] (0,-0.6) rectangle (0.6,0.6);
  \end{scope}
  \draw[thick] (0,0) circle (0.55);
  \node[font=\footnotesize] at (0,0) {地球};
  \foreach \y in {-2.4,-1.2,0,1.2,2.4} {
    \draw[-Stealth, omThm, thick] (7.7,\y) -- (6.4,\y);
  }
  \node[right, font=\small, omThm] at (7.8,0.55) {阳光};
  \node[left, font=\footnotesize] at (-4.85,0) {满月};
  \node[right, font=\footnotesize] at (4.85,0) {新月};
  \node[above, font=\footnotesize] at (0,3.75) {下弦月};
  \node[below, font=\footnotesize] at (0,-3.75) {上弦月};
\end{tikzpicture}

{\small 一张图看懂全部奥秘：绕行中的\omterm{def:g2:moon-in-the-sky:moon}{月球}八个位置，
每个位置朝向太阳的一侧都被照亮。变化的只是
向阳面中哪一部分朝向地球。}
\end{omfigure}

\begin{method}[做地球，托住月球]\label{met:g7:moon-phases:walk}
下面的解释，一分钟即可演示：
\begin{enumerate}
  \item 在暗室中，一盏裸灯代表太阳；你代表
        地球；你伸直的手上托着的白球代表
        \omterm{def:g2:moon-in-the-sky:moon}{月球}；
  \item 面向灯，让白球朝向灯：白球的受光半球背向你，
        你看到的是它的暗面：新月；
  \item 原地慢慢转身（白球绕着你公转）：球的
        边缘出现一道亮边——\omterm{def:g7:moon-phases:phases}{渐盈}的\omterm{ex:g2:moon-in-the-sky:shapes}{蛾眉月}；
        转过四分之一圈，一半的球被照亮：
        上弦月；
  \item 继续转动：凸月——然后，灯在身后，白球
        正对你的一面全部被照亮：\omterm{ex:g2:moon-in-the-sky:shapes}{满月}；再转完整的一圈，
        经过\omterm{def:g7:moon-phases:phases}{渐亏}的半月与\omterm{ex:g2:moon-in-the-sky:shapes}{蛾眉月}回到新月。
\end{enumerate}
你自转一圈，便重演了一个月的天象。
\end{method}

\section{已被摒弃的误解}

\begin{proposition}[月相不是地球的影子]\label{prop:g7:moon-phases:notshadow}
天文学中最常见的错误答案是：“\omterm{def:g7:moon-phases:phases}{月相}就是
地球在\omterm{def:g2:moon-in-the-sky:moon}{月球}上的\omterm{def:g1:light-and-shadows:shadow}{影子}。”几何学从三个方面驳斥了它。
地球的\omterm{def:g1:light-and-shadows:shadow}{影子}\emph{背离}太阳伸展——它只能
触及位于太阳对面的\omterm{def:g2:moon-in-the-sky:moon}{月球}，这只在
\omterm{ex:g2:moon-in-the-sky:shapes}{满月}时发生（而且只有在罕见的\omterm{def:g7:shadows-eclipses:lunar}{月食}中才会）。\omterm{ex:g2:moon-in-the-sky:shapes}{蛾眉月}
\omterm{def:g2:moon-in-the-sky:moon}{月球}位于我们的\emph{向日}一侧，远离我们的\omterm{def:g1:light-and-shadows:shadow}{影子}。而且
\omterm{def:g1:light-and-shadows:shadow}{影子}在\omterm{def:g2:moon-in-the-sky:moon}{月球}上的边缘总是一个大圆的一部分，
——它永远无法刻画出凸月那种双重弯曲的
边界。\omterm{def:g7:moon-phases:phases}{月相}无需\omterm{def:g1:light-and-shadows:shadow}{影子}：一个半被照亮的球，从
不同的方位看去，便能解释一切。
\end{proposition}

\begin{example}[月食与月相]\label{ex:g7:moon-phases:eclipsevs}
要把两种月亮变暗的现象分开存放。\emph{\omterm{def:g7:moon-phases:phases}{月相}}：
即缓慢的月度循环，\omterm{def:g2:moon-in-the-sky:moon}{月球}自身的背阴面转向
我们——与我们的\omterm{def:g1:light-and-shadows:shadow}{影子}无关。\emph{\omterm{def:g7:shadows-eclipses:lunar}{月食}}：在一个罕见的
\omterm{ex:g2:measuring-time:clock}{小时}里，\omterm{ex:g2:moon-in-the-sky:shapes}{满月}穿过我们的\omterm{prop:g7:shadows-eclipses:umbra}{本影}——地球的\omterm{def:g1:light-and-shadows:shadow}{影子}
真实地投在它身上，呈铜红色，来去迅速。月亮日记中的节律是
一种视角，\omterm{def:g7:shadows-eclipses:lunar}{月食}则是一次造访。
\end{example}

\section{像钟表一样读懂天空}

\begin{example}[何时月相保有其时刻]\label{ex:g7:moon-phases:hours}
图示将每个\omterm{def:g7:moon-phases:phases}{月相}固定在一个位置上，
而每个位置又对应一张时间表。\omterm{ex:g2:moon-in-the-sky:shapes}{满月}位于太阳的对面：它应在\omterm{def:g1:day-and-night:daynight}{日落}时升起，
整夜高挂，并在\omterm{def:g1:day-and-night:daynight}{日出}时落下。新月
\emph{与}太阳同行：整日升起，却淹没在强光之中——
这正是没有人“看见”新月的原因。上弦月落后太阳四分之一圈，
它大约在中午升起，午夜前后落下——即晚饭后的月亮。你童年时对白天之月的惊奇，
如今已成为一张时刻表：
凡既非新月也非\omterm{ex:g2:moon-in-the-sky:shapes}{满月}的月亮，
都与太阳共享一部分的运行时段。
\end{example}

\begin{example}[蛾眉月的尖角指向何方？]\label{ex:g7:moon-phases:points}
任一\omterm{def:g7:moon-phases:phases}{月相}的亮面都朝向太阳——\omterm{ex:g2:moon-in-the-sky:shapes}{蛾眉月}像一张对准射手的弓。
傍晚时分，\omterm{ex:g2:moon-in-the-sky:shapes}{蛾眉月}位于西方：两端的尖角向上、向外伸展，
背离地平线下方落日的余晖。
天空从不会画错，画家却常常画错。尖角朝向太阳的\omterm{ex:g2:moon-in-the-sky:shapes}{蛾眉月}，
是天文学中最爱考的“找错”题。
\end{example}

\begin{remark}[那张永不转开的脸]\label{rem:g7:moon-phases:sameface}
日记中还有一条观察，终于得到了解释：\omterm{def:g2:moon-in-the-sky:moon}{月球}上的暗色斑块
每晚、每个\omterm{def:g7:moon-phases:phases}{月相}、年复一年，总是呈现同一面。\omterm{def:g2:moon-in-the-sky:moon}{月球}绕自身的轴
自转恰好一周，正好完成一次公转，因此总有同一半面朝向地球。
直到一个空间探测器携带相机绕到背面之前，没有人见过\omterm{def:g2:moon-in-the-sky:moon}{月球}的背面，
而这发生在你祖父母还记得的年代。
转动与公转为何如此完美地吻合，
是一个关于潮汐的深刻故事，
留待本课程的大学阶段再讲。
\end{remark}

\section{练习}

\begin{exercise}[$\star$]\label{exo:g7:moon-phases:1}
任一时刻，\omterm{def:g2:moon-in-the-sky:moon}{月球}被太阳照亮的部分占多大比例？
那么，什么是“\omterm{def:g7:moon-phases:phases}{月相}”？
\end{exercise}

\begin{exercise}[$\star$]\label{exo:g7:moon-phases:2}
从新月开始，按顺序背出各个\omterm{def:g7:moon-phases:phases}{月相}，并给出周期的长度。
周期中哪一半是“\omterm{def:g7:moon-phases:phases}{渐盈}”？
\end{exercise}

\begin{exercise}[$\star$]\label{exo:g7:moon-phases:3}
在沿\omterm{def:g6:sun-earth-moon:axisorbit}{轨道}行走的方法中，新月时灯、球和你各在何处？
\omterm{ex:g2:moon-in-the-sky:shapes}{满月}时呢？
\end{exercise}

\begin{exercise}[$\star$]\label{exo:g7:moon-phases:4}
给出三个几何论据中的两个，用以反驳
“\omterm{def:g7:moon-phases:phases}{月相}是地球的\omterm{def:g1:light-and-shadows:shadow}{影子}”这一说法。
\end{exercise}

\begin{exercise}[$\star$]\label{exo:g7:moon-phases:5}
新月的黑暗，
与\omterm{def:g7:shadows-eclipses:lunar}{月食}时的黑暗有何区别？
\end{exercise}

\begin{exercise}[$\star$]\label{exo:g7:moon-phases:6}
为什么\omterm{ex:g2:moon-in-the-sky:shapes}{满月}总是在接近\omterm{def:g1:day-and-night:daynight}{日落}时升起？
请依据图示中它的位置加以论证。
\end{exercise}

\begin{exercise}[$\star$]\label{exo:g7:moon-phases:7}
为什么没有人能“看见新月在午夜升起”？
午夜时，新月位于何处？
\end{exercise}

\begin{exercise}[$\star$]\label{exo:g7:moon-phases:8}
傍晚时分，一弯\omterm{ex:g2:moon-in-the-sky:shapes}{蛾眉月}挂在西方。
相对于已经落下的太阳，它的尖角指向何方——
为什么\omterm{ex:g2:moon-in-the-sky:shapes}{蛾眉月}的尖角永远不可能指向太阳？
\end{exercise}

\begin{exercise}[$\star\star$]\label{exo:g7:moon-phases:9}
今晚\omterm{def:g2:moon-in-the-sky:moon}{月球}正处于\omterm{def:g7:moon-phases:phases}{渐盈}凸月阶段。它在\omterm{def:g6:sun-earth-moon:axisorbit}{轨道}上位于何处？
它大约何时升起？五个晚上之后，会是什么\omterm{def:g7:moon-phases:phases}{月相}，
又遵循怎样的时间表？
\end{exercise}

\begin{exercise}[$\star\star$]\label{exo:g7:moon-phases:10}
\omterm{def:g7:shadows-eclipses:solar}{日食}只能发生在某一种\omterm{def:g7:moon-phases:phases}{月相}，\omterm{def:g7:shadows-eclipses:lunar}{月食}只能发生在另一种\omterm{def:g7:moon-phases:phases}{月相}。
请指出这两种\omterm{def:g7:moon-phases:phases}{月相}，并依据三者排成一线的几何关系加以说明，
解释为什么在这两种\omterm{def:g7:moon-phases:phases}{月相}时，
\omterm{def:g7:shadows-eclipses:solar}{日食}与\omterm{def:g7:shadows-eclipses:lunar}{月食}仍然很少发生。
\end{exercise}

\begin{exercise}[$\star\star$]\label{exo:g7:moon-phases:11}
位于\omterm{def:g2:moon-in-the-sky:moon}{月球}正面的宇航员仰望地球。
描述一个月内他们看到的“地球\omterm{def:g7:moon-phases:phases}{月相}”，
并说明当我们看到新月时，地球呈现给他们的是什么\omterm{def:g7:moon-phases:phases}{月相}。
（两个方向的几何关系都要算。）
\end{exercise}

\begin{exercise}[$\star\star\star$]\label{exo:g7:moon-phases:12}
旧月的秘密：晴朗的傍晚，\omterm{ex:g2:moon-in-the-sky:shapes}{蛾眉月}
\omterm{def:g2:moon-in-the-sky:moon}{月球}的\emph{暗色}部分会微微发光——“新月臂弯中的旧月”。
你二年级时曾把它当作谜题来接触；
现在请彻底解开它：追踪光的路径，
解释为什么地照在恰为\omterm{ex:g2:moon-in-the-sky:shapes}{蛾眉月}的\omterm{def:g7:moon-phases:phases}{月相}时最亮，
此时\emph{地球}向\omterm{def:g2:moon-in-the-sky:moon}{月球}呈现的是什么\omterm{def:g7:moon-phases:phases}{月相}？
以及为什么它在接近\omterm{ex:g2:moon-in-the-sky:shapes}{满月}时逐渐变暗。
\end{exercise}

\section{问题：与月亮共度一月}

\begin{problem}[{周末问题——天文台的开放月；三十夜的月亮，
一张挂图；讲解员的
问答}]\label{pb:g7:moon-phases:1}
镇上的天文台开展一项公开活动，名为“与月亮共度一月”。
作为初级讲解员，你负责保管挂图，并回答参观者的问题。

\medskip\noindent\textbf{第一部分——设置挂图。}
\begin{enumerate}
  \item 第 0 夜为新月。填写挂图上的关键\omterm{def:g5:series-parallel-circuits:parallel}{节点}：大约在哪些夜晚
        出现上弦月、\omterm{ex:g2:moon-in-the-sky:shapes}{满月}和
        下弦月？
  \item 挂图要求在每个关键\omterm{def:g5:series-parallel-circuits:parallel}{节点}旁标注升起时刻。
        配对：\omterm{def:g1:day-and-night:daynight}{日落}时升起；\ 正午前后升起；\ 午夜前后
        升起；\ 与太阳同升同落。
  \item 一位参观者问：挂图上的\omterm{ex:g2:moon-in-the-sky:shapes}{蛾眉月}低挂在
        \emph{西}方的傍晚天空，而\omterm{ex:g2:moon-in-the-sky:shapes}{满月}却在
        黄昏时从\emph{东}方升起。请结合
        \omterm{def:g6:sun-earth-moon:axisorbit}{轨道}位置加以解释。
  \item 另一位参观者坚称自己在上午十点
        看到了月亮。挂图上哪些\omterm{def:g7:moon-phases:phases}{月相}使
        清晨看到月亮成为可能？
\end{enumerate}

\medskip\noindent\textbf{第二部分——讲解员的
演示。}
\begin{enumerate}[resume]
  \item 你为一个学生团体演示“灯-球行走”实验。
        一名学生在你转身时喊停：球体显示
        四分之三被照亮。请说出\omterm{def:g7:moon-phases:phases}{月相}，并为团体指出灯、
        你在地球上的位置以及球体。
  \item 学生指着你那暗面的球问：“所以是地球的\omterm{def:g1:light-and-shadows:shadow}{影子}造成\omterm{def:g7:moon-phases:phases}{月相}？”
        请给出演示本身的反驳：他们看到的
        是谁的夜半球？
        地球的\omterm{def:g1:light-and-shadows:shadow}{影子}实际上
        应落在何处？
  \item 团体成员问：为什么总能在月面上看到“玉兔”（或“人脸”），
        却从来看不到它的背面？请用
        “同一面”的事实作答。
  \item 一位持怀疑态度的家长问：“如果\omterm{def:g2:moon-in-the-sky:moon}{月球}像你说的那样在自转，
        我们不应该看到它转动吗？”请解释这一表面上的矛盾：
        \omterm{def:g2:moon-in-the-sky:moon}{月球}怎样能在每次绕地球公转一周时恰好自转一圈，
        从而\emph{隐藏}了自身的自转？
\end{enumerate}

\medskip\noindent\textbf{第三部分——问答之夜。}
\begin{enumerate}[resume]
  \item 问答卡 A：“今晚的月亮在午夜升起，
        其\emph{左}半边被照亮（从朝向
        清晨的方向看去）。\omterm{def:g7:moon-phases:phases}{渐盈}还是\omterm{def:g7:moon-phases:phases}{渐亏}？属于上弦还是下弦？”
  \item 问答卡 B：“一幅艺术海报画着午夜的\omterm{ex:g2:moon-in-the-sky:shapes}{蛾眉月}，
        高挂天顶，月角笔直指向
        地平线上\omterm{def:g1:day-and-night:daynight}{日落}时的色彩。”请列出
        海报上的每一处天文错误。
  \item 问答卡 C：“今晚\omterm{ex:g2:moon-in-the-sky:shapes}{满月}时，参观者能否
        同时在地球某处目睹\omterm{def:g7:shadows-eclipses:solar}{日食}？”
        请用\omterm{def:g7:moon-phases:phases}{月相}的几何关系加以判断。
  \item 结束致辞：用三句话说明本月的一个基本事实，
        本月所掩盖的一个误解，以及
        参观者应带回家的时刻表窍门
        （任一\omterm{def:g7:moon-phases:phases}{月相}都说明了它的升起时刻）。
\end{enumerate}
\end{problem}
